\section{Introduction}
We collect proposed answers, partial results and prior-work deductions
concerning the September 2026 revision of the 21st edition of the
Kourovka Notebook~\cite{kourovka21}. The index below records the answered
subparts and leads directly to the arguments. Problem statements are
concise restatements of the Notebook, with the original authors credited;
the hypotheses and convention-dependent limits of each answer remain
explicit. References to other solutions distinguish shared arguments,
different constructions and stronger results without assigning discovery
priority.

These arguments originated in a 48-hour experiment organized by Michael
Kinyon and Josef Urban, using OpenAI's GPT-6 Astra through Codex, GAP and
other exact computations. Anthropic's Claude Opus[1m], through Claude Code,
subsequently reviewed the work. The full research-and-review account,
\emph{Forty-eight hours with the Kourovka Notebook}, gives the experimental
conditions, reports, correspondence and revision history in the same
\href{https://github.com/JUrban/kour1}{public repository}.
The present edition uses the same mathematical source files.

The original deadline list contained 46 candidate entries. The claimed
answer to 10.35 was withdrawn after Evgeny Khukhro pointed out that its
target field is $\overline{\Q}$, not $\Q$: the proposed example already
embeds over $\Q(i)$. Both model reviews had missed this error. The
withdrawn application is omitted here and documented in the full account.
The remaining 45 entries are proposed answers, not a certified count of
correct new theorems. In particular, the cases with explicit convention
limits and the older ingredients identified for 16.9 retain that status.

The arguments are ordered by problem number. Partial results follow in
Section~\ref{app:partials}; rediscoveries, deductions and a formulation
issue are separated in Section~\ref{app:prior}. Computational proof inputs
and their availability are specified in Appendix~\ref{app:math-artifacts}.
Throughout, $[x,y]=x^{-1}y^{-1}xy$.
